\begin{lemma}[Tangential bound and the Laplacian in the adapted frame]
Let $n,m\in\mathbb{N}$ with $n\ge3$ and $m\ge2$, let
$P\in\mathbb{R}[X_1,\dots,X_n]$ be homogeneous of degree $m$, let $\mathbb{R}^n$
carry the Euclidean inner product and let $u(y)=P(y)$. Put $\theta_0=\pi/m$. Let
$e,v_0\in\mathbb{R}^n$ satisfy $\lVert e\rVert=\lVert v_0\rVert=1$ and
$\langle e,v_0\rangle=0$, and assume that $v_0$ maximizes the function
$v\mapsto u(\cos(\theta_0)e+\sin(\theta_0)v)$ over all unit vectors $v$ with
$\langle e,v\rangle=0$, i.e.
\[
  u\bigl(\cos(\theta_0)e+\sin(\theta_0)v\bigr)
   \le u\bigl(\cos(\theta_0)e+\sin(\theta_0)v_0\bigr)
  \qquad\text{for every such }v .
\]
Write $F(s)=u(\cos(s)e+\sin(s)v_0)$ and
$x_0=\cos(\theta_0)e+\sin(\theta_0)v_0$. Then for every unit vector $w$ with
$\langle e,w\rangle=\langle v_0,w\rangle=0$,
\[
  \bigl\langle\nabla u(x_0),w\bigr\rangle=0,
  \qquad
  D^2u(x_0)[w,w]\le m\,F(\theta_0)+F'(\theta_0)\,
    \frac{\cos\theta_0}{\sin\theta_0},
\]
and moreover
\[
  \Delta u(x_0)\ \le\ m(m+n-2)\,F(\theta_0)+F''(\theta_0)
    +(n-2)\,F'(\theta_0)\,\frac{\cos\theta_0}{\sin\theta_0}.
\]
\end{lemma}
\begin{proof}
Put $e_r=x_0=\cos(\theta_0)e+\sin(\theta_0)v_0$ and
$e_\theta=-\sin(\theta_0)e+\cos(\theta_0)v_0$.

\emph{Step 0: elementary facts.} Since $m\ge2$ and $\pi>0$ we have
$0<\theta_0=\pi/m\le\pi/2$, so $\sin\theta_0>0$ ($\sin$ is strictly positive on
$(0,\pi)$). Expanding inner products and using
$\lVert e\rVert=\lVert v_0\rVert=1$, $\langle e,v_0\rangle=0$:
\begin{equation}\label{eq:tlb-frame}
  \lVert e_r\rVert^2=\cos^2\theta_0+\sin^2\theta_0=1,
  \quad
  \lVert e_\theta\rVert^2=\sin^2\theta_0+\cos^2\theta_0=1,
  \quad
  \langle e_r,e_\theta\rangle=-\cos\theta_0\sin\theta_0+\sin\theta_0\cos\theta_0=0,
\end{equation}
so $(e_r,e_\theta)$ is an orthonormal pair. Also, directly,
\begin{equation}\label{eq:tlb-decomp}
  \sin(\theta_0)\,e_r+\cos(\theta_0)\,e_\theta=v_0,
  \qquad
  \cos(\theta_0)\,e_r-\sin(\theta_0)\,e_\theta=e ,
\end{equation}
as one checks by expanding and using $\cos^2\theta_0+\sin^2\theta_0=1$.

\emph{Step 1: the comparison curves.} Fix a unit vector $w$ with
$\langle e,w\rangle=\langle v_0,w\rangle=0$ and set, for $\sigma\in\mathbb{R}$,
\[
  v_\sigma=\cos(\sigma)v_0+\sin(\sigma)w,
  \qquad
  \eta(\sigma)=\cos(\theta_0)e+\sin(\theta_0)\,v_\sigma .
\]
Expanding and using that $v_0,w$ are unit and mutually orthogonal, and that both
are orthogonal to $e$,
\[
  \lVert v_\sigma\rVert^2=\cos^2\sigma+\sin^2\sigma=1,
  \qquad
  \langle e,v_\sigma\rangle=\cos(\sigma)\langle e,v_0\rangle+\sin(\sigma)\langle e,w\rangle=0 .
\]
So each $v_\sigma$ is a unit vector orthogonal to $e$. Note also
$v_0$ corresponds to $\sigma=0$, hence $\eta(0)=x_0$. Finally, writing
\[
  C=\cos(\theta_0)\,e,\qquad A=\sin(\theta_0)\,v_0,\qquad B=\sin(\theta_0)\,w,
\]
we have $\eta(\sigma)=C+\cos(\sigma)A+\sin(\sigma)B$, which is exactly the curve
shape of \noderef{GreatCircleSecondDeriv}; its derivative at $\sigma$ is
$-\sin(\sigma)A+\cos(\sigma)B$, so at $\sigma=0$ it equals
$B=\sin(\theta_0)w$.

\emph{Step 2: $h(\sigma)=u(\eta(\sigma))$ has a maximum at $\sigma=0$.} By
Step 1 each $v_\sigma$ is a unit vector orthogonal to $e$, so the maximality
hypothesis applies to it and gives
\[
  h(\sigma)=u\bigl(\cos(\theta_0)e+\sin(\theta_0)v_\sigma\bigr)
  \le u\bigl(\cos(\theta_0)e+\sin(\theta_0)v_0\bigr)=u(x_0)=h(0).
\]

\emph{Step 3: the first-order condition.} By \noderef{MvPolyFDeriv} applied to
$P$ at the point $\eta(0)=x_0$, the polynomial function $z\mapsto P(z)$ on the
coordinate space is Frechet differentiable at $x_0$ with derivative
$z\mapsto\sum_i(\partial_iP)(x_0)z_i$; composing with the identity map of
$\mathbb{R}^n$ read from the Euclidean to the coordinate norm, which is a
continuous linear isomorphism because any two norms on a finite-dimensional
space are equivalent, the chain rule makes $u$ Frechet differentiable at $x_0$
with $Du(x_0)(z)=\sum_i(\partial_iP)(x_0)z_i$, and by \noderef{MvPolyGradient}
together with the coordinate formula for the inner product this equals
$\langle\nabla u(x_0),z\rangle$. The curve $\eta$ is differentiable (its
coordinates are constants plus linear combinations of $\cos$ and $\sin$), so by
the chain rule $h$ is differentiable at $0$ with
\[
  h'(0)=Du(x_0)\bigl(\sin(\theta_0)w\bigr)=\sin(\theta_0)\,\langle\nabla u(x_0),w\rangle .
\]
By Step 2 the differentiable function $h$ has an interior maximum at $0$, so
$h'(0)=0$; dividing by $\sin\theta_0>0$ gives
\begin{equation}\label{eq:tlb-first}
  \bigl\langle\nabla u(x_0),w\rangle=0 ,
\end{equation}
which is the first assertion.

\emph{Step 4: the second-order condition.} By \noderef{GreatCircleSecondDeriv}
applied with the $C,A,B$ of Step 1 at the parameter $\sigma=0$, and using
$\eta(0)=x_0$ and $\cos(0)A+\sin(0)B=A=\sin(\theta_0)v_0$,
\[
  h''(0)=D^2u(x_0)\bigl[\sin(\theta_0)w,\sin(\theta_0)w\bigr]
   -\bigl\langle\nabla u(x_0),\,\sin(\theta_0)v_0\bigr\rangle .
\]
By \noderef{MvPolyHessian}, $D^2u(x_0)[z,z']=\sum_i\sum_j z_i z'_j(\partial_i\partial_jP)(x_0)$
is linear in each argument separately, so the first term equals
$\sin^2(\theta_0)\,D^2u(x_0)[w,w]$; and the inner product is linear in its second
argument, so
\begin{equation}\label{eq:tlb-second}
  h''(0)=\sin^2(\theta_0)\,D^2u(x_0)[w,w]
   -\sin(\theta_0)\,\langle\nabla u(x_0),v_0\rangle .
\end{equation}
The function $h$ is in fact twice differentiable at $0$. Indeed, by
\noderef{MvPolyContDiff} the function $u$ is of class $C^k$ on $\mathbb{R}^n$ for
every smoothness exponent $k$, and the curve $\eta$ is infinitely
differentiable, since each of its coordinates is a constant plus a linear
combination of $\cos$ and $\sin$ and a map into a finite-dimensional space is
$C^k$ exactly when its coordinates are; so $h=u\circ\eta$ is infinitely
differentiable by the chain rule. (This is what makes $h''(0)$ in
\eqref{eq:tlb-second} the genuine second derivative of $h$ at $0$.) By Step 2
the function $h$ has a maximum at the interior point $0$, and a function that is
twice differentiable at an interior maximum has nonpositive second derivative
there. Hence $h''(0)\le0$, and \eqref{eq:tlb-second} gives
\[
  \sin^2(\theta_0)\,D^2u(x_0)[w,w]\le\sin(\theta_0)\,\langle\nabla u(x_0),v_0\rangle .
\]
Dividing by $\sin^2(\theta_0)>0$,
\begin{equation}\label{eq:tlb-div}
  D^2u(x_0)[w,w]\le\frac{\langle\nabla u(x_0),v_0\rangle}{\sin\theta_0}.
\end{equation}

\emph{Step 5: evaluating $\langle\nabla u(x_0),v_0\rangle$.} By
\eqref{eq:tlb-decomp} and linearity of the inner product in its second
argument,
\[
  \langle\nabla u(x_0),v_0\rangle
  =\sin(\theta_0)\,\langle\nabla u(x_0),e_r\rangle
   +\cos(\theta_0)\,\langle\nabla u(x_0),e_\theta\rangle .
\]
By \noderef{FrameHessian} applied with this $e$, this $v_0$ and $t=\theta_0$
(its hypotheses are $m\ge1$, which follows from $m\ge2$, and the homogeneity of
$P$), the first two of its five identities give
$\langle\nabla u(x_0),e_r\rangle=m\,F(\theta_0)$ and
$\langle\nabla u(x_0),e_\theta\rangle=F'(\theta_0)$. Hence
\[
  \langle\nabla u(x_0),v_0\rangle
  =\sin(\theta_0)\,m\,F(\theta_0)+\cos(\theta_0)\,F'(\theta_0),
\]
and substituting into \eqref{eq:tlb-div} and dividing through by
$\sin\theta_0>0$,
\begin{equation}\label{eq:tlb-tang}
  D^2u(x_0)[w,w]\ \le\ m\,F(\theta_0)
    +F'(\theta_0)\,\frac{\cos\theta_0}{\sin\theta_0},
\end{equation}
which is the second assertion.

\emph{Step 6: the adapted orthonormal basis.} By \eqref{eq:tlb-frame} the pair
$(e_r,e_\theta)$ is an orthonormal family of two vectors in $\mathbb{R}^n$, and
$n\ge3\ge2$, so by the background proposition on extending an orthonormal family
there is an orthonormal basis $(b_i)_{i=1}^{n}$ of $\mathbb{R}^n$ and two
distinct indices $i_0\ne i_1$ with $b_{i_0}=e_r$ and $b_{i_1}=e_\theta$. Let
$i\notin\{i_0,i_1\}$. Then $b_i$ is a unit vector with
$\langle e_r,b_i\rangle=\langle e_\theta,b_i\rangle=0$, so by
\eqref{eq:tlb-decomp} and linearity,
\[
  \langle e,b_i\rangle=\cos(\theta_0)\langle e_r,b_i\rangle-\sin(\theta_0)\langle e_\theta,b_i\rangle=0,
  \qquad
  \langle v_0,b_i\rangle=\sin(\theta_0)\langle e_r,b_i\rangle+\cos(\theta_0)\langle e_\theta,b_i\rangle=0 .
\]
So each such $b_i$ is a unit vector orthogonal to both $e$ and $v_0$, i.e. an
admissible choice of $w$ in \eqref{eq:tlb-tang}.

\emph{Step 7: summing the frame.} By the background proposition computing the
Laplacian from an orthonormal basis, applied with $(b_i)_{i=1}^{n}$,
\[
  \Delta u(x_0)=\sum_{i=1}^{n}D^2u(x_0)[b_i,b_i]
  =D^2u(x_0)[e_r,e_r]+D^2u(x_0)[e_\theta,e_\theta]
   +\sum_{i\notin\{i_0,i_1\}}D^2u(x_0)[b_i,b_i],
\]
splitting off the two indices $i_0\ne i_1$ and using $b_{i_0}=e_r$,
$b_{i_1}=e_\theta$. By \noderef{FrameHessian} at $t=\theta_0$, the third and
fifth of its identities give
\[
  D^2u(x_0)[e_r,e_r]=m(m-1)F(\theta_0),
  \qquad
  D^2u(x_0)[e_\theta,e_\theta]=F''(\theta_0)+m\,F(\theta_0).
\]
By Step 6 and \eqref{eq:tlb-tang}, each of the remaining $n-2$ summands (the
index set $\{i:i\notin\{i_0,i_1\}\}$ has exactly $n-2$ elements) is at most
$m\,F(\theta_0)+F'(\theta_0)\cos\theta_0/\sin\theta_0$, so their sum is at most
$(n-2)$ times that. Combining,
\[
  \Delta u(x_0)\le m(m-1)F(\theta_0)+F''(\theta_0)+m\,F(\theta_0)
    +(n-2)\Bigl(m\,F(\theta_0)+F'(\theta_0)\frac{\cos\theta_0}{\sin\theta_0}\Bigr).
\]
Collecting the coefficients of $F(\theta_0)$, namely
$m(m-1)+m+(n-2)m=m\bigl((m-1)+1+(n-2)\bigr)=m(m+n-2)$, this is exactly
\[
  \Delta u(x_0)\le m(m+n-2)F(\theta_0)+F''(\theta_0)
    +(n-2)F'(\theta_0)\frac{\cos\theta_0}{\sin\theta_0},
\]
the third assertion.
\end{proof}
